\chapter{Geluid: snelheid, toonhoogte, luidheid}\label{ch:g8:sound-pitch-loudness}

Korte snaren zingen hoog, harde tokkels zingen luid --- je
elastiekjesgitaar vond de regels jaren geleden al. Dit jaar krijgen de
regels getallen: een telling per \omterm{def:g2:measuring-time:second}{seconde} voor de toonhoogte, een
zwaaigrootte voor de luidheid, en een venster --- met randen --- op wat
menselijke oren überhaupt kunnen opvangen. Voorbij de ene rand lachen
de vleermuizen ons uit.

\section{Frequentie: toonhoogte geteld}

\begin{definition}[Frequentie]\label{def:g8:sound-pitch-loudness:frequency}
De \emph{frequentie}\index{frequentie} $f$ van een
\omterm{def:g3:sound-around-us:vibration}{trilling} is het aantal volledige heen-en-weerzwaaien dat ze per
\omterm{def:g2:measuring-time:second}{seconde} maakt. Haar \omterm{def:g6:measurement-in-science:measuring}{eenheid} is de \emph{hertz}\index{hertz}
(\unit{Hz}): één zwaai per \omterm{def:g2:measuring-time:second}{seconde} --- ter ere van de natuurkundige die
als eerste onzichtbare golven maakte en opving. Een gitaarsnaar die
$110$ keer per \omterm{def:g2:measuring-time:second}{seconde} trilt, trilt op \qty{110}{Hz}; de vleugels van
een mug op bijna \qty{600}{Hz}; het trage klapperen van een grote vlag
op enkele hertz.
\end{definition}

\begin{proposition}[Toonhoogte is frequentie]\label{prop:g8:sound-pitch-loudness:pitch}
Wat het oor \emph{toonhoogte} noemt, is
\omterm{def:g8:sound-pitch-loudness:frequency}{frequentie} en niets anders: hoe sneller de bron trilt, hoe
hoger de noot. De elastiekjesregels vertalen zich meteen --- korte,
strakke, dunne dingen trillen op hogere $f$. Het anker van de muziek:
concertinstrumenten over de hele wereld stemmen op de noot a van
\qty{440}{Hz}. Het verdubbelen van de \omterm{def:g8:sound-pitch-loudness:frequency}{frequentie} van welke noot
dan ook tilt haar met het welluidendste interval van allemaal --- het
octaaf: weer een a, een octaaf hoger, op \qty{880}{Hz}.
\end{proposition}

\begin{omfigure}
\begin{tikzpicture}
\begin{axis}[omaxis, width=10.8cm, height=4.4cm,
    xlabel={tijd (\unit{ms})}, ylabel={zwaai},
    xtick={0,5,10,15,20}, ytick=\empty,
    ymin=-1.5, ymax=1.5, xmin=0, xmax=21, samples=200]
  \addplot[very thick, omDef, domain=0:20]
    {sin(deg(2*3.14159*0.1*x))};
  \addplot[very thick, omThm, domain=0:20]
    {0.9*sin(deg(2*3.14159*0.4*x))};
\end{axis}
\end{tikzpicture}

{\small Twee \omterm{def:g3:sound-around-us:sound}{geluiden} getekend door een microfoon: in dezelfde $20$
milliseconden zwaait het blauwe spoor twee keer (\qty{100}{Hz}, een lage
bromtoon) terwijl het rode acht keer zwaait (\qty{400}{Hz}, een gezongen
noot). Toonhoogte is de telling.}
\end{omfigure}

\begin{method}[Geluid zien]\label{met:g8:sound-pitch-loudness:seeing}
De trommel met rijstkorrels toonde het trillen van het
\omterm{def:g3:sound-around-us:sound}{geluid}; een microfoon tekent het.
\begin{enumerate}
  \item draai een geluidsanalyseprogramma op een telefoon of computer
        --- het scherm zet de zwaai van de microfoon uit tegen de
        tijd, een moderne oscilloscoop;
  \item neurie laag en gelijkmatig: een luie, wijd uiteen liggende
        lijn; zing hoog: de pieken dringen samen;
  \item lees de \omterm{def:g8:sound-pitch-loudness:frequency}{frequentie} af terwijl je fluit: de meeste
        fluittonen landen tussen $500$ en \qty{2000}{Hz};
  \item sla de stemvork met het opschrift ``a 440'' aan en controleer
        of de aflezing zijn gravure eer aandoet.
\end{enumerate}
Eén scherm vervangt een hoofdstuk giswerk: toonhoogte geteld voor je
ogen.
\end{method}

\section{Het venster van het oor}

\begin{proposition}[Het bereik van het menselijk gehoor]\label{prop:g8:sound-pitch-loudness:range}
Gezonde jonge oren vangen \omterm{def:g3:sound-around-us:vibration}{trillingen} op tussen ongeveer
\qty{20}{Hz} en \qty{20000}{Hz}. Onder het venster ligt
\emph{infrageluid}\index{infrageluid} --- soms als gerommel in de borst
gevoeld, nooit gehoord; erboven \emph{ultrageluid}\index{ultrageluid}
--- stilte voor ons, hoe heftig het trillen ook is. Het venster wordt
smaller met de leeftijd en met misbruik: het hoogste octaaf is
doorgaans het eerste slachtoffer van luide gewoontes.
\end{proposition}

\begin{example}[Voorbij de randen]\label{ex:g8:sound-pitch-loudness:beyond}
De natuur en de industrie werken allebei voorbij ons venster.
Vleermuizen jagen met ultrasone klikjes rond \qty{50000}{Hz} en lezen
\omterm{ex:g4:how-sound-travels:echo}{echo}'s fijner dan enig hoorbaar \omterm{def:g3:sound-around-us:sound}{geluid} zou kunnen
schilderen --- dolfijnen net zo, in de geest van het sonarhoofdstuk.
Honden antwoorden op ``stille'' fluitjes die net boven onze rand
liggen. Medische scanners \omterm{prop:g5:mirrors-reflection:image}{beelden} ongeboren kinderen af met
onschadelijke ultrageluidsecho's. Aan de andere rand converseren
olifanten in infrasoon gerommel over kilometers, en stormen en
aardbevingen kondigen zich aan in tonen waar geen oor naar luistert.
\end{example}

\section{Luidheid: de grootte van de zwaai}

\begin{proposition}[Luidheid is amplitude]\label{prop:g8:sound-pitch-loudness:loudness}
Wat het oor \emph{luidheid} noemt, komt overeen met de
\emph{amplitude}\index{amplitude} van de \omterm{def:g3:sound-around-us:vibration}{trilling} --- de grootte
van de zwaai, de sterkte van de reizende samenpersing. Tokkel zachtjes
of woest: de pieken van de lijn staan laag of hoog terwijl hun
\emph{onderlinge afstand} --- de toonhoogte --- onaangeroerd blijft.
\omterm{prop:g8:sound-pitch-loudness:loudness}{Amplitude} en \omterm{def:g8:sound-pitch-loudness:frequency}{frequentie} zijn onafhankelijke knoppen:
elke toonhoogte mag gefluisterd of gebulderd worden.
\end{proposition}

\begin{omfigure}
\begin{tikzpicture}
\begin{axis}[omaxis, width=10.8cm, height=4.4cm,
    xlabel={tijd (\unit{ms})}, ylabel={zwaai},
    xtick={0,5,10,15,20}, ytick=\empty,
    ymin=-1.5, ymax=1.5, xmin=0, xmax=21, samples=200]
  \addplot[very thick, omDef, domain=0:20]
    {0.35*sin(deg(2*3.14159*0.25*x))};
  \addplot[very thick, omThm, domain=0:20]
    {1.2*sin(deg(2*3.14159*0.25*x))};
\end{axis}
\end{tikzpicture}

{\small Dezelfde noot, gefluisterd en gebulderd: gelijke afstand tussen
de pieken (één toonhoogte, \qty{250}{Hz}), ongelijke piekhoogte ---
luidheid is de grootte van de zwaai.}
\end{omfigure}

\begin{example}[De decibelladder]\label{ex:g8:sound-pitch-loudness:decibels}
Luidheid wordt gegradeerd in \emph{decibel}\index{decibel} (\unit{dB})
--- een ladder gebouwd voor het verbluffende bereik van het oor, waarop
elke vertienvoudiging van de sterkte van het \omterm{def:g3:sound-around-us:sound}{geluid} hetzelfde
aantal sporten toevoegt. IJkpunten: ritselende bladeren \qty{10}{dB};
een rustig gesprek \qty{50}{dB}; druk verkeer \qty{80}{dB}; de eerste
rijen van een rockconcert \qty{110}{dB}; pijn en onmiddellijke schade
rond \qty{120}{dB}. De vreemde rekenkunde van de ladder --- gelijke
stappen voor tienvoudige sterktes --- is een logaritme, een van de
grote uitvindingen die je op de middelbare school opwachten.
\end{example}

\begin{remark}[Het budget van het oor]\label{rem:g8:sound-pitch-loudness:earsafety}
Het trommeltje in je oor, dat je in je hoofdstuk uit groep vijf
ontmoette, houdt levenslang de boekhouding bij. Boven ongeveer
\qty{85}{dB} loopt de schade op met de blootstellingsduur --- een
werkdag als limiet in een lawaaiige fabriek, slechts minuten op het
niveau van de eerste concertrijen; een koptelefoon op vol volume zit
midden in de gevarenband. Gehoor dat aan de hoge zwaaien verloren
gaat, komt nooit terug: de hoogste tonen vertrekken het eerst, geruisloos
en voorgoed. Oren hebben geen oogleden --- en geen reserveonderdelen.
\end{remark}

\begin{example}[Snelheid, onberoerd door beide knoppen]\label{ex:g8:sound-pitch-loudness:speed}
Geen van beide knoppen raakt de bezorging aan. Hoog en laag, luid en
zacht, alle \omterm{def:g3:sound-around-us:sound}{geluiden} rijden mee op de ene estafette van het
\omterm{def:g7:sound-production:medium}{medium} en op het ene tempo van dat \omterm{def:g7:sound-production:medium}{medium} ---
\qty{340}{m/s} in \omterm{def:g2:air-around-us:air}{lucht}. Het bewijs wordt elke avond opgevoerd:
een orkest dat aan de overkant van een park wordt gehoord, komt
\emph{op toon en op tijd} aan --- waren hoge noten sneller dan lage, dan
zou elk ver akkoord tot arpeggio's uitsmeren. (De luidheid vervaagt met
de afstand --- de uitdijende schillen uit het hoofdstuk van vorig jaar
--- maar vervagen is verzwakken, nooit vertragen.)
\end{example}

\section{Oefeningen}

\begin{exercise}[$\star$]\label{exo:g8:sound-pitch-loudness:1}
Definieer \omterm{def:g8:sound-pitch-loudness:frequency}{frequentie} en haar \omterm{def:g6:measurement-in-science:measuring}{eenheid}. Wat is de
toonhoogte van de concert-a, en van de a een octaaf hoger?
\end{exercise}

\begin{exercise}[$\star$]\label{exo:g8:sound-pitch-loudness:2}
Twee lijnen beslaan dezelfde $20$ milliseconden: de ene toont $4$
volledige zwaaien, de andere $16$. Bereken beide \omterm{def:g8:sound-pitch-loudness:frequency}{frequenties}; welke
klinkt hoger?
\end{exercise}

\begin{exercise}[$\star$]\label{exo:g8:sound-pitch-loudness:3}
Geef het venster van het menselijk gehoor, en de namen van de twee
landen buiten zijn grenzen --- met van elk één inwoner.
\end{exercise}

\begin{exercise}[$\star$]\label{exo:g8:sound-pitch-loudness:4}
Wat verandert er op het scherm van de analyser wanneer je dezelfde noot
luider zingt? En wanneer je een hogere noot even luid zingt?
\end{exercise}

\begin{exercise}[$\star$]\label{exo:g8:sound-pitch-loudness:5}
Vertaal de oude elastiekjesregels naar frequentietaal: wat doen korter,
strakker en dunner elk met $f$?
\end{exercise}

\begin{exercise}[$\star$]\label{exo:g8:sound-pitch-loudness:6}
Plaats op de decibelladder: ritselende bladeren; verkeer; de eerste rij
van het concert; pijn. Waar begint de gevarenband?
\end{exercise}

\begin{exercise}[$\star$]\label{exo:g8:sound-pitch-loudness:7}
Waarom lijken hondenfluitjes voor hun baasjes kapot? En waarom is de
hond het daar niet mee eens?
\end{exercise}

\begin{exercise}[$\star$]\label{exo:g8:sound-pitch-loudness:8}
Een ver orkest komt op toon en op tijd aan. Wat bewijst dit over de
\omterm{def:g7:motion-average-speed:formula}{snelheid} van \omterm{def:g3:sound-around-us:sound}{geluid} over toonhoogte en luidheid
heen?
\end{exercise}

\begin{exercise}[$\star\star$]\label{exo:g8:sound-pitch-loudness:9}
Een vleermuis klikt op \qty{50000}{Hz}. Hoeveel volledige zwaaien
bevat één klik van \qty{0.01}{s}? Waarom moet fijne
echobeeldvorming zo snel trillen? (Denk aan wat de fijnste
penselen gemeen hebben.)
\end{exercise}

\begin{exercise}[$\star\star$]\label{exo:g8:sound-pitch-loudness:10}
De laagste toets van een vleugel trilt rond \qty{27}{Hz}; de hoogste
rond \qty{4200}{Hz}. Plaats beide binnen (of tegenover) het venster van
het gehoor --- en leg uit waarom een verouderende pianostemmer aan
één kant van het klavier het eerst de hulp van de analyser nodig kan
hebben.
\end{exercise}

\begin{exercise}[$\star\star$]\label{exo:g8:sound-pitch-loudness:11}
Concertregels staan $8$ uur toe bij \qty{85}{dB} maar slechts
minuten rond \qty{110}{dB}, waarbij het budget om de paar sporten
halveert. Welke soort groei verbergt zich in de gelijke stappen van de
ladder --- en wat leert het instorten van het budget over ``nog een
klein beetje luider''?
\end{exercise}

\begin{exercise}[$\star\star\star$]\label{exo:g8:sound-pitch-loudness:12}
Ontwerp de volledige geluidscheck van een schoolaula met één
analysertelefoon: hoe je de a van de piano tegen $440$ verifieert; hoe
je de luidheid van de zaal op de achterste rij tijdens een drumtest
vindt; hoe je betrapt of de ventilatie een infrasoon gerommel uitzendt
dat het oor mist; en welk van de drie bevindingen de microfoon van de
telefoon het \emph{minst} kan bevestigen (raadpleeg het eigen venster
van het gehoor van de telefoon --- gebouwd voor stemmen).
\end{exercise}

\section{Probleem: de geluidscheck}

\begin{problem}[{Weekendprobleem --- geluidscheck op het
schoolfestival; één analyser, drie bands en een zenuwachtige
verpleegkundige; de avond van de mengtafel}]\label{pb:g8:sound-pitch-loudness:1}
Vooravond van het festival: jij bedient de analyser aan de mengtafel,
met de schoolverpleegkundige die over je schouder naar de decibelmeter
kijkt.

\medskip\noindent\textbf{Deel I --- Het stemuur.}
\begin{enumerate}
  \item De stemvork van het koor zingt op \qty{440}{Hz}; de analyser
        toont de a van de piano op \qty{432}{Hz}. Stel de diagnose van
        de piano, en noem de knop --- toonhoogte of luidheid --- die de
        stemmer nodig heeft.
  \item De laagste snaar van de basgitaar leest \qty{41}{Hz}. Plaats
        hem in het venster van het gehoor --- en leg uit waarom het
        publiek hem net zo goed in de borst zal \emph{voelen} als
        horen.
  \item De aangehouden noot van de leadzanger toont pieken met precies
        dezelfde onderlinge afstand als die van de stemvork, maar drie
        keer zo hoog. Vergelijk de twee \omterm{def:g3:sound-around-us:sound}{geluiden} op beide
        knoppen.
  \item Een jonge sjouwer beweert dat de nieuwe speakers van de zaal
        ``tot \qty{22000}{Hz} weergeven''. Wie in het publiek zal de
        bovenkant van dat bereik kunnen bevestigen --- en wie niet?
\end{enumerate}

\medskip\noindent\textbf{Deel II --- Het grootboek van de
verpleegkundige.}
\begin{enumerate}[resume]
  \item Repetitieniveaus: akoestische set \qty{75}{dB}; rockset
        \qty{100}{dB}; finale gepland ``een tikje hoger'' op
        \qty{110}{dB}. Welke sets markeert de verpleegkundige, en tegen
        welke budgetregel?
  \item De verpleegkundige deelt schuimoordoppen uit die
        \omterm{def:g3:sound-around-us:sound}{geluiden} met \qty{20}{dB} dempen. Wat wordt de rockset
        erachter --- en waarom dragen verstandige sjouwers bij elke
        show doppen?
  \item Op tweemaal zo grote afstand van het podium zakt de decibelmeter
        merkbaar. Welke oude wet van de uitdijende schillen is aan het
        werk --- en welke knop (toonhoogte of luidheid) draait de
        afstand?
  \item De drummer protesteert: ``Zachter betekent doffer --- de hoge
        noten sterven het eerst als we zacht spelen!'' Ontwar de twee
        knoppen voor hem: wat verandert zacht spelen werkelijk?
\end{enumerate}

\medskip\noindent\textbf{Deel III --- Festivalnatuurkunde.}
\begin{enumerate}[resume]
  \item Vuurwerk sluit de show af op \qty{680}{m} afstand. Bereken de
        vertraging tussen flits en knal die het publiek zal merken ---
        en de \omterm{def:g8:sound-pitch-loudness:frequency}{frequentie} van het gedreun als de grote bommen twee
        keer per \omterm{def:g2:measuring-time:second}{seconde} dreunen (venstercontrole: gehoord of
        gevoeld?).
  \item Het laatste mysterie van de avond: een telefoonvideo van het
        concert klinkt bij het terugkijken dun --- het gerommel in
        de borst is weg. Raadpleeg het voorbehoud van oefening 12: wat
        was de microfoon nooit gebouwd om op te vangen?
  \item Het verslag van de verpleegkundige vraagt om ``elk één zin''
        over de twee knoppen van de avond en het derde ding dat geen
        van beide knoppen kan aanraken. Stel ze op.
  \item Afbreken: de ultrasone ongediertverjager van de conciërge, die
        de hele avond bij de podiumdeur draaide, stoorde geen enkele
        luisteraar en, naar verluidt, veel muizen. Verklaar beide
        publieken met één venstervoorstel.
\end{enumerate}
\end{problem}
